% Discussion. Derived from manuscript/discussion.argdown after Round 1
% reviewer integration (LaTeX prose).

% --- Headline ---
In this pre-registered 16-trial pilot, the strict-tolerance EFDPR of 0.31 (Clopper-Pearson 95\% CI 0.11--0.59) exceeded the pre-registered 0.25 threshold as a point estimate, but the pre-registered one-sided exact-binomial test did not reject the null at $\alpha = 0.05$ ($P = 0.37$); the corresponding liberal-tolerance estimate sits exactly at the threshold. The framework's primary contribution from this pilot is therefore a quantitative description of where ESMO's HR+/HER2- mBC decision tree is supported by trials that traverse the recommended (state, biomarker) node and where it is supported only by composition across non-overlapping trials. The decision points identified as evidence-free under strict concordance --- post-CDK4/6i ESR1mut, AKT-pathway, PIK3CAmut, no-actionable-mutation, and the post-CDK4/6i+post-endo everolimus node --- are precisely the choices clinicians face daily and precisely the positions in the guideline tree where pivotal-trial enrolment criteria most often differ from the guideline-node patient state.

\paragraph{Pilot scale and the gap between estimate and inference.}
The 16-trial / 16-node corpus is the pre-registered pilot, not a systematic-review-scale census. With $n=16$, the one-sided exact-binomial test of $H_0: p \le 0.25$ requires $k \ge 8$ successes to reject at $\alpha = 0.05$; the observed $k=5$ does not. The gap between a point estimate that exceeds the threshold and a test that does not reject is a property of the pilot's sample size, not of the framework. A production-scale extension that adds the remaining HR+/HER2- mBC pivotal trials registered between 2015 and 2026, together with their guideline-citation network captured via OpenAlex, will resolve this in the planned methods paper.

\paragraph{Composition-only citations are routine but should be visible.}
Guideline committees routinely compose recommendations across non-overlapping trials when no single trial covers the exact decision node. ESMO's citation of SOLAR-1 for the post-CDK4/6i PIK3CAmut node (when SOLAR-1 enrolled a post-AI population), and of CAPItello-291 for the post-CDK4/6i AKT-pathway node (when CAPItello-291 enrolled a mixed pre/post-CDK4/6i population with only a subgroup readout), are canonical examples. The framework does not declare these compositions illegitimate; it makes them explicit and counts them. Whether a given composition is clinically acceptable is a separate, downstream question that benefits from being able to see, on a single diagram, where composition has been invoked and which alternative trial chains could have been considered.

\paragraph{Algorithm transparency over clinical reasoning.}
The strict concordance algorithm is deliberately conservative: it over-flags reasonable clinical extrapolations as evidence-free so that the resulting count is a robust upper bound on directly-supported decisions rather than a clinical judgement about extrapolation. The three pre-registered tolerance levels (strict / ESCAT-aligned / liberal) are the framework's mechanism for quantifying the dependence of EFDPR on extrapolation assumptions: in this pilot, moving from strict to liberal converted only one node (G5, post-CDK4/6i ESR1mut) from evidence-free to supported, indicating that the gap is not a tolerance artefact.

\paragraph{Reproducible by construction.}
Every analytic choice in this study is materialised as either a JSON file or a Python function and released under MIT licence at the repository linked in the Code Availability statement. The 16-trial structured-eligibility set is Supplementary Table S2, the 16-node ESMO encoding is Supplementary Table S3, and the biomarker-token vocabulary for the ODI computation is Supplementary Table S4. The schema is frozen at v1.0; the tagged release \texttt{v1.0.0} reproduces every figure and table.

\paragraph{Extensible across tumour types.}
The (prior-state, biomarker) node space is tumour-agnostic; tumour-specific biomarker tokens enter as additive vocabulary. The concordance algorithm is biomarker-agnostic. The same pipeline is planned for application to EGFR-mutant non-small-cell lung cancer and to metastatic colorectal cancer with ctDNA-guided decision axes; in each case the only artefacts that need authoring are the structured trial corpus, the structured guideline decision tree, and the per-tumour ESCAT grouping table.

\paragraph{Actionable for guideline committees and trial sponsors.}
A committee that adopts this framework as a guideline-update audit can locate the specific decision nodes that lack supporting evidence chains. Trial sponsors can use the list of evidence-free post-CDK4/6i biomarker-stratified nodes to design enrichment trials whose target population matches a guideline-recommended decision point, rather than the convenience population most readily recruited from a single network. The EFDPR overlay adds one column to an artefact (the decision-tree table) that committees already produce.

\paragraph{Limitations.}
Three limitations bound the present work. First, the 16-trial corpus is a pilot drawn from the ESMO/ASCO reference lists, which introduces selection toward trials the guidelines already cite; the production-grade extension will replace this seed with a systematic ClinicalTrials.gov query. Second, eligibility coding was hand-curated from canonical primary publications; the LLM-extraction pipeline planned for production has its own validation requirements (per-field accuracy and Cohen's $\kappa$ against this hand-curated set) and is reported separately. Third, the primary analysis uses ESMO 2024 as the single guideline source; ASCO and NCCN sensitivity analyses are pre-registered for the production extension. The internal multi-round review process identified six NCT-to-trial-name misidentifications across two rounds (full correction log in Supplementary Text 1), which were resolved against ClinicalTrials.gov \texttt{briefTitle} and \texttt{conditionsModule} fields before the present draft was finalised; we report this transparency in the spirit of pre-registered, audit-friendly reporting.

\paragraph{Outlook.}
The framework, schema, and ESMO encoding are released so that guideline committees, trial sponsors, and meta-research groups can extend and audit them. A community-maintained \emph{Treatment-Sequencing Evidence Atlas}, refreshed at every ESMO and ASCO guideline release, could provide a continuously-updated EFDPR overlay across solid tumours; immediate priorities are the NSCLC and mCRC applications, the LLM-extraction validation paper with inter-rater agreement against the hand-curated set, and an interactive web viewer that lets a clinician click any guideline decision node and see, in one screen, every trial-DAG edge that supports it and every edge that does not.
